% Section 6.4: projected gain of each design class at the operating point
% (Qwen2.5-72B fp8, 16-bit KV, GH200, ctx 14724, B=3; r chosen by
% eq. decision for each design, 2026-10-22). Ben, 2026-10-21:
% "projected graph comparison with DirectKV/Pie/DAK". Values printed by
% paper/model/compare_designs.py; each is the capacity model applied to the
% mechanism the paper describes, not a model of the system.
\begin{figure}[t]
\centering
\begin{tikzpicture}
\begin{axis}[
  width=0.72\columnwidth, height=4.2cm,
  xbar, bar width=5pt,
  xmin=0.8, xmax=2.0,
  ymin=0.4, ymax=4.6,
  ytick={1,2,3,4},
  yticklabels={HBM only, aggregation (DAK), {zero-copy (DirectKV)},
               \sysname{}},
  xtick={0.8,1.0,1.2,1.4,1.6,1.8},
  xticklabel style={/pgf/number format/.cd,fixed,fixed zerofill,precision=1},
  xlabel={projected throughput gain over HBM only},
  tick label style={font=\sffamily\scriptsize},
  label style={font=\sffamily\scriptsize},
  axis line style={black!60}, tick style={black!60},
  xmajorgrids, grid style={black!10},
  legend style={font=\sffamily\scriptsize, draw=none, fill=white,
                fill opacity=0.85, text opacity=1, at={(0.5,-0.32)},
                anchor=north, legend columns=2, column sep=6pt},
  legend cell align=left,
  nodes near coords={\pgfmathprintnumber[fixed,fixed zerofill,precision=2]{\pgfplotspointmeta}},
  nodes near coords style={font=\sffamily\tiny, black!80},
  point meta=x,
]
\addplot[draw=hbmD, fill=hbmD!35, bar shift=3pt]
  coordinates {(1.000,1) (1.105,2) (0.897,3) (1.499,4)};
\addlegendentry{mechanism only}
\addplot[draw=ddrD, fill=ddrD, bar shift=-3pt]
  coordinates {(1.368,4)};
\addlegendentry{with \sysname{}'s measured overheads}
\draw[black!60, densely dashed, line width=0.6pt]
  (axis cs:1.794,0.4) -- (axis cs:1.794,4.6);
\node[font=\sffamily\tiny, black!70, anchor=south west, align=left]
  at (axis cs:1.81,1.6) {link\\ceiling,\\no costs\\(1.79)};
\end{axis}
\end{tikzpicture}
\caption{\textbf{Projected gain by design}, at the operating point of
\S\ref{sec:opoint} on a GH200 (context 14\,724, three residents). Each bar is
the capacity model applied to the mechanism each paper describes, not a model
of that system. Aggregation reads the same batch faster and is bounded by
$1+\Xlink/BW$. Zero-copy of every request's KV is bounded by the link, $\Xlink/(S\btok)$
requests per unit time. \sysname{} adds requests over the link, with the
rider count of \cref{eq:decision}. The light bar omits implementation
overheads; the dark bar adds \sysname{}'s measured $c_{\mathrm{m}}$ and
$F$.}
\label{fig:compare}
\end{figure}
